% Pipeline overview figure for the BioTest paper.
% Spans both columns of the IEEEtran conference template.
% Two-row block diagram in the minimal academic style.
%
% Visual semantics (no embedded legend; caption does the work):
%   - rounded rectangle, light fill, blue rule: framework phase
%   - rounded rectangle, white fill, grey rule: data artefact
%   - rounded rectangle, darker fill, heavier rule: user-supplied or decision step
%   - solid arrow: forward data flow
%   - dashed arrow: Phase D coverage feedback into Phase B

\tikzset{
  bt/box/.style={
    rectangle, rounded corners=1.5pt,
    draw=black!55, fill=white,
    line width=0.4pt,
    inner sep=2.5pt, minimum height=7mm,
    align=center, font=\scriptsize
  },
  bt/phase/.style={
    bt/box,
    draw=blue!50!black, fill=blue!4,
    line width=0.5pt
  },
  bt/oracle/.style={
    bt/box,
    draw=blue!55!black, fill=blue!10,
    line width=0.7pt
  },
  bt/user/.style={
    bt/box,
    draw=black!80, fill=black!8,
    line width=0.6pt
  },
  bt/flow/.style={
    -{Latex[length=1.3mm,width=1.1mm]},
    line width=0.45pt, draw=black!75
  },
  bt/feedback/.style={
    -{Latex[length=1.3mm,width=1.1mm]},
    dashed, line width=0.45pt, draw=black!55
  },
}

\begin{figure*}[t]
\centering
\begin{tikzpicture}[font=\scriptsize, node distance=3.5mm and 3.5mm]

% --- Row 1: one-shot, per-format ----------------------------------
\node[bt/box, text width=14mm] (spec)
  {VCF / SAM\\spec};
\node[bt/phase, right=of spec, text width=22mm] (A)
  {\textbf{Phase A}\\Ingest \& embed\\\textit{(ChromaDB)}};
\node[bt/phase, right=of A, text width=30mm] (B)
  {\textbf{Phase B}\\LLM-RAG MR mining\\\textit{deepseek-v4-flash}};
\node[bt/box, right=of B, text width=23mm] (mr)
  {MR registry\\\textit{seed-validated}};
\node[bt/box, right=of mr, text width=34mm] (stack)
  {Corpus stack\\Ranks 1--7 metamorphic\\Ranks 8--13 \textbf{(Phase E)}\\augmentation};

\draw[bt/flow] (spec) -- (A);
\draw[bt/flow] (A) -- (B);
\draw[bt/flow] (B) -- (mr);
\draw[bt/flow] (mr) -- (stack);

% --- Row 2: per-iteration -----------------------------------------
\node[bt/user, below=11mm of A.south, text width=22mm] (shim)
  {\textbf{User shim}\\\textit{parse} $\to$ JSON};
\node[bt/phase, right=of shim, text width=24mm] (C)
  {\textbf{Phase C}\\Multi-runner exec\\$N{\geq}3$ SUTs};
\node[bt/oracle, right=of C, text width=27mm] (oracle)
  {\textbf{Consensus oracle}\\differential $+$ metamorphic};
\node[bt/box, right=of oracle, text width=28mm] (out)
  {Candidate bugs\\$+$ mutation grading};

\draw[bt/flow] (shim) -- (C);
\draw[bt/flow] (C) -- (oracle);
\draw[bt/flow] (oracle) -- (out);

% Corpus stack feeds Phase C (drop between rows, lower channel)
\draw[bt/flow]
  (stack.south) |- ($(C.north)+(0,1.5mm)$) -- (C.north);

% Phase D coverage feedback: oracle output back into Phase B
% Routed in the upper channel between rows so it does not overlap
% the corpus-drop arrow.
\draw[bt/feedback]
  (oracle.north) -- ++(0,6mm) -|
  node[pos=0.18, above, font=\scriptsize, text=black!65]
       {\textit{Phase D: coverage steering}}
  (B.south);

\end{tikzpicture}
\caption{\biotest{} pipeline. Phases~A and~B run once per format,
and phases~C, D, and~E run per iteration. The dashed arrow marks
Phase~D coverage feedback into Phase~B. The user's only per-SUT
artefact is the parse-and-emit shim, the same artefact in-process
coverage-guided fuzzers require.}
\label{fig:pipeline}
\end{figure*}
